\begin{frame}{Thuật toán BFS (Breadth-First Search)}
    \begin{itemize}
        \item \textbf{Nguyên lý:} Duyệt theo từng lớp, bắt đầu từ đỉnh nguồn, lan rộng ra các đỉnh láng giềng.
        \item \textbf{Cấu trúc dữ liệu:} Sử dụng \textbf{Hàng đợi (Queue)} theo cơ chế FIFO (Vào trước - Ra trước).
        \item \textbf{Các bước:}
        \begin{enumerate}
            \item Đưa đỉnh nguồn vào Queue và đánh dấu đã thăm.
            \item Lặp khi Queue không rỗng:
            \begin{itemize}
                \item Lấy đỉnh $u$ ra khỏi Queue.
                \item Xét các đỉnh $v$ kề với $u$: Nếu chưa thăm, đưa $v$ vào Queue và đánh dấu.
            \end{itemize}
        \end{enumerate}
    \end{itemize}
\end{frame}

% Slide 2: Code C++
\begin{frame}[fragile]{Cài đặt BFS với C++}
\begin{verbatim}
void BFS(int start) {
    queue<int> q;
    q.push(start);
    visited[start] = true;
    while(!q.empty()){
        int u = q.front(); q.pop();
        cout << u << " ";
        for(int v : adj[u]){
            if(!visited[v]){
                visited[v] = true;
                q.push(v);
            }
        }
    }
}
\end{verbatim}
\end{frame}